\section{P\&A y temas de debate}

\subsection{Resumen de preguntas y respuestas del público}

En la segunda mitad de la entrevista, Burak respondió preguntas que le interesaban al público:

\begin{itemize}
    \item \textbf{P: ¿Cómo lograr que el agent se comporte de manera consistente en contextos multilingües?}\\
    R: Construir bien el pipeline de localización, internacionalizar tanto el prompt como el tool manifest, y usar una pila de embedding multilingüe.
    \item \textbf{P: ¿Cómo auditar la hallucination del agent?}\\
    R: Combinar search+LLM: al generar la answer, producir también la source cite correspondiente, y enviar tanto la cite como el user prompt al evaluation harness.
\end{itemize}

\subsection{Temas de debate y reflexión}

Burak planteó dos proposiciones que vale la pena debatir:

\begin{enumerate}
    \item ¿El guardrail de seguridad del agent debe ir primero o iterarse de manera simultánea?
    \item ¿Las capacidades Multi-modal deben desplegarse primero en task de bajo riesgo?
\end{enumerate}

Su postura es la siguiente: el guardrail necesita iterarse de manera simultánea (no se puede esperar a que la capacidad mejore para agregarlo); lo Multi-modal debe validarse primero en escenarios poco sensibles a la high-bandwidth observation, y luego extenderse gradualmente a escenarios de alta sensibilidad.

\subsection{Resumen de la sección}

La parte de P\&A reveló el pain point que más le interesaba al público, y también puso de relieve que las empresas deben mantener un debate constante entre la capacidad y la governance.

